\chapter{Introdução}

O plano global criado pela ONU, denominado Agenda 2030, propõe uma ação conjunta que visa o desenvolvimento sustentável e a erradicação da pobreza em todas as suas formas. Segundo \cite{onu2015agenda2030}, a ação proposta pela Agenda 2030 estabelece 17 objetivos e 169 metas que abrangem as cinco grandes áreas: pessoas, planeta, prosperidade, paz e parceria. As cinco áreas abrangem objetivos para garantir condições dignas de vida às pessoas, possibilitando acesso à educação, saúde e igualdade. Conservação do meio ambiente, com o objetivo de combater a degradação ambiental e promover uso sustentável de recursos ambientais, possibilitando um desenvolvimento e prosperidade da sociedade de forma equilibrada, visando também estabelecer a paz e a justiça social. 

O Brasil possui o Índice de Desenvolvimento Sustentável das Cidades (IDSC) \cite{idscbr2023}, 100 indicadores coletados que apresentam métricas de cada uma das cidades com relação ao cumprimento dos 17 objetivos de desenvolvimento sustentável (ODS) definidos na Agenda. Esses indicadores oferecem uma visão do desenvolvimento dos municípios de forma multidimensional, compreendendo esferas sociais, ambientais e institucionais. Um dos indicadores pertencentes ao IDSC é o Índice de Desenvolvimento da Educação Básica (IDEB), responsável por avaliar a qualidade da educação básica. De acordo com \cite{inep_ideb}, o indicador, calculado a partir de dados, abrange dois aspectos fundamentais: o fluxo escolar e as médias de desempenho nas avaliações. 

Avanços tecnológicos, como Modelos de Aprendizado de Máquina, fornecem técnicas de aprendizado supervisionado para a predição de avaliadores como o IDEB. A predição desse valor é realizada exclusivamente com base nos dados do IDSC, diferente da realização do cálculo baseado em notas e taxas de aprovação. Sendo assim, leva em conta a realidade do aluno com base em fatores externos à sala de aula. A abordagem aplicada fornece uma predição do valor que se baseia em aspectos sociais, econômicos e ambientais do município dos alunos. A abordagem é válida para gestão pública educacional, apresentando o valor do IDEB integrado a fatores da realidade dos alunos fora do ambiente escolar. 

A predição do valor de avaliação do desempenho educacional ainda é aplicada em conjunto com modelos de aprendizado não supervisionados, que são responsáveis por agrupar o conjunto de dados com base em características. Dessa forma, os dados do IDSC podem ser agrupados com relação à um conjunto predefinido de variáveis. O agrupamento do conjunto de dados se mostra uma tecnologia essencial nesse cenário, visto que o conjunto de dados do IDSC, em sua totalidade, abrange 5.570 municípios de diferentes perfis de cidades brasileiras. Os indicadores de desenvolvimento sustentável podem apresentar relações distintas com o índice do IDEB, dependendo das características do município.

Além da segmentação do conjunto de dados para a modelagem do indicador de avaliação do IDEB, se faz necessária a interpretabilidade do resultado da predição. Para a resolução desse aspecto, serão aplicadas métricas de relevância dos atributos nos modelos baseados em árvores de decisão, o que permite extrair de forma direta quais os indicadores da base de dados do IDSC que possuem maior influência no índice educacional. Embora a análise desses fatores seja baseada principalmente nas métricas extraídas dos modelos, ainda é utilizada Inteligência Artificial Explicável (XAI) como forma de comparação entre ambas as ferramentas destinadas à interpretabilidade do modelo. Diante do cenário apresentado, o entendimento dos aspectos que têm maior peso na variação do índice do IDEB auxiliam à tomada de decisão no poder público, visando à melhoria do sistema educacional. A tarefa de identificação dos fatores de desenvolvimento sustentável do município que influenciam diretamente a qualidade da educação básica pode levar à implementação de políticas públicas com objetivos diretos na educação. 

\section{Problema}

O conjunto de dados IDSC, que abrange todos os indicadores de cumprimento dos objetivos definidos na Agenda 2030, apresenta diversas características com relação ao desenvolvimento sustentável do município que ainda não foram analisadas e exploradas sob a ótica da ciência de dados sobretudo de forma aprofundada por meio da construção de modelos capazes de realizar a predição do indicador do IDEB.

O conjunto abrange informações sobre todos os municípios brasileiros, mas as características de cada localidade podem ter diferentes níveis de associação com o IDEB, não apresentando uma relação determinante com o indicador de avaliação em todos os municípios de forma igualitária.

O indicador do IDEB geralmente é visto como uma métrica isolada apenas no âmbito da educação, sem levar em conta aspectos do desenvolvimento do município, como saúde ou saneamento básico. Aspectos esses que estão disponíveis no conjunto de indicadores IDSC e ainda não foram explorados com esse objetivo. Os modelos preditivos podem ser aplicados e, com base no treinamento, predizer a nota do IDEB. Porém, para uma tomada de decisão com relação à melhoria do sistema educacional, se faz necessária a compreensão dos fatores que influenciam o resultado da predição, por isso a proposta visa à aplicabilidade de técnicas de interpretabilidade dos modelos empregados. 

\section{Objetivos}

Este trabalho se propõe à exploração e análise da base de dados dos indicadores do Índice de Desenvolvimento Sustentável das Cidades (IDSC), aplicando modelos de aprendizado de máquina com o objetivo de realizar o agrupamento com base em características selecionadas do município e predição do Índice de Desenvolvimento da Educação Básica (IDEB).

Adicionalmente, busca-se analisar as métricas de relevância de atributos inerentes aos modelos preditivos para interpretar os fatores que mais influenciam os resultados. O objetivo geral deste trabalho é analisar o desempenho de modelos de regressão baseados em Árvores de Decisão, comparando sua aplicação na base de dados completa e em conjuntos agrupados por características dos município, utilizando suas métricas estruturais de importância de atributos para identificar os indicadores sustentáveis de maior impacto na educação básica.

Para alcançar o objetivo geral, definem-se os seguintes objetivos específicos:

 Realizar uma revisão da literatura sobre modelos de aprendizado de máquina supervisionado e não supervisionado, com ênfase em modelos baseados em árvores de decisão e métricas de interpretabilidade de atributos;
 Aplicar técnicas para exploração e análise inicial da base de dados do IDSC;
 Aplicar técnicas de modelo de aprendizado não supervisionado para agrupamento do conjunto de dados com base no perfil socioeconômico e sustentável do município;
 Aplicar técnicas de modelo de aprendizado de máquina de regressão baseados em árvores de decisão para prever o valor do avaliador educacional IDEB;
 Aplicar métricas de extração da importância dos atributos nos modelos preditivos, aplicar também abordagens de IA Explicável (XAI) de forma complementar para validar a influência dos indicadores do IDSC no IDEB e o entendimento do funcionamento do modelo de predição;
 Avaliar e comparar o desempenho dos modelos e o comportamento das métricas de relevância de atributos entre a base integral e as bases agrupadas.


\section{Justificativa}

Diante da crescente complexidade do cenário social e econômico presente nos municípios brasileiros, a tomada de decisão baseada em evidências se revela essencial para a formulação de políticas públicas. O Brasil apresenta acentuada heterogeneidade entre seus municípios devido à grande extensão territorial e à diversidade cultural e socioeconômica. Dessa forma, a adoção de políticas públicas uniformes tende a ser ineficaz, pois falha em abranger todos os cenários e desafios de cada município brasileiro. 

Nesse cenário, técnicas de análise de dados e aprendizado de máquina possibilitam a identificação de padrões complexos entre os indicadores de desenvolvimento sustentável (IDSC) e o Índice de Desenvolvimento da Educação Básica (IDEB). O agrupamento dos municípios conforme características semelhantes permite uma análise refinada e contextualizada conforme a realidade de cada perfil regional. A partir desse agrupamento, a predição do índice fornece uma visão estratégica fundamental aos gestores públicos, permitindo identificar possíveis relações entre fatores sociais e ambientais e o rendimento escolar, o que possibilita a adoção de medidas prévias para a melhoria da educação básica, antes da consolidação do índice oficial.

Para que essa aplicação seja verdadeiramente útil à gestão pública, a predição do modelo deve vir acompanhada da interpretabilidade. A predição com a aplicação de modelos baseados em árvores de decisão permite extrair diretamente a relevância de atributos. Embora amparada pelos conceitos de Inteligência Artificial Explicável (XAI), a análise se concentra no uso dessas métricas para mapear quais fatores do IDSC exercem maior peso sobre o desempenho educacional. Esse diagnóstico quantitativo direciona o poder público para a implementação de políticas em aspectos específicos, solucionando problemas que são a causa principal da desigualdade escolar.

\section{Metodologia}

Para a realização de todos os objetivos, será necessária uma exploração inicial da base de dados, além de um estudo aprofundado em técnicas de aprendizado supervisionado, não supervisionado e métodos de interpretabilidade de modelos. Essa exploração da base e o estudo teórico visam compreender o cenário e fundamentar a escolha dos algoritmos de aprendizado de máquina, levando em consideração uma separação por agrupamento consistente dos dados e melhores resultados na predição do índice do IDEB. A análise exploratória de dados será realizada com o objetivo de compreender a distribuição dos indicadores e identificar padrões preliminares. Após essa análise, será executado o pré-processamento da base de dados, visando à adequação dos dados para a execução dos modelos de aprendizado de máquina, de forma a garantir a confiabilidade e o desempenho dos algoritmos.

A etapa preditiva inicial consiste na aplicação dos modelos de regressão na base de dados integral para a obtenção de um resultado preditivo de referência. Os modelos aplicados para comparação e análise serão Regressão Linear, para base inicial de comparação por se caracterizar como um modelo mais simples em comparação com os demais, e os algoritmos baseados em árvores de decisão Random Forest e Gradient Boosting, sendo representado neste trabalho pelos algoritmos \textit{Extreme Gradient Boosting}, \textit{Light Gradient Boosting Machine} e \textit{Categorical Boosting}. Logo após, para a identificação de perfis de municípios semelhantes, serão aplicadas técnicas de aprendizado não supervisionado. Nessa etapa de agrupamento, serão analisados e avaliados os algoritmos \textit{K-Means}, Agrupamento Hierárquico e Agrupamento Espacial Baseado em Densidade para Aplicações com Ruído.

A etapa subsequente consiste na realização da predição em cada subconjunto da base de dados, definidos pelo modelo de agrupamento que apresentar o melhor desempenho. Como estratégia de aprendizado supervisionado, serão replicados os mesmos modelos aplicados na base integral, permitindo avaliar e comparar o impacto do agrupamento dos municípios em bases de dados menores na precisão dos resultados. O desempenho dos modelos de predição será avaliado por meio das métricas de Erro Médio Absoluto ($MAE$), Raiz do Erro Quadrático Médio ($RMSE$) e Coeficiente de Determinação ($R^2$). 

Por fim, a extração da relevância dos atributos ocorrerá por meio de duas abordagens complementares. A primeira consiste na análise das métricas de importância intrínsecas aos algoritmos de árvore, calculadas a partir da redução da variância dos dados nos nós de decisão. Em seguida, para mitigar os vieses estruturais dessas métricas e garantir maior consistência teórica, será adotado o método \textit{TreeSHAP}. Enquanto as métricas nativas oferecem uma perspectiva estrutural de quais indicadores do IDSC foram mais acionados pelos modelos, os valores SHAP quantificam a magnitude e direção do impacto de cada variável. Essa abordagem permite identificar se um indicador atua de forma positiva ou negativa na estimativa do IDEB.